% !TeX program = lualatex
% !TeX root = predavanje02.tex
% ============================================================
% Zajednički preamble za sva predavanja RWA kolegija
% Uključiti s: % ============================================================
% Zajednički preamble za sva predavanja RWA kolegija
% Uključiti s: % ============================================================
% Zajednički preamble za sva predavanja RWA kolegija
% Uključiti s: \input{../common/preamble.tex}
% ============================================================

\documentclass[aspectratio=149,11pt]{beamer}
\usetheme{CambridgeUS}

\usepackage{fontspec}              % radi samo sa XeLaTeX/LuaLaTeX
\setsansfont{Latin Modern Sans}    % sans-serif font
\setmonofont{Latin Modern Mono}    % monospace font za kod
\usepackage[croatian]{babel}
\usepackage{csquotes}              % \enquote{} → hrvatski navodnici „..."
\usepackage{lmodern}
\usepackage{booktabs}
\usepackage{listings}
\usepackage{xcolor}
\usepackage{hyperref}

% ---------- Autor i logo ----------
\author{mag.ing. Josip Torić (SIT UNIZD)}
\titlegraphic{\includegraphics[height=2.5cm]{../common/SIT_logo_tekst1.pdf}}

% ---------- Listings ----------
\definecolor{codebg}{RGB}{248,248,248}
\definecolor{codeframeclr}{RGB}{220,220,220}
\lstdefinelanguage{bash}{
  morekeywords={sudo,apt,brew,pip,pipx,uv,winget,python,git,flask,source,deactivate,code,gunicorn,docker,compose,heroku,render,railway,node,npm,pnpm,yarn,uvicorn,alembic,pytest,ruff,black,precommit},
  sensitive=false,
  morecomment=[l]{\#}
}
\lstdefinelanguage{python}{
  morekeywords={from,import,def,class,return,if,elif,else,for,in,while,try,except,raise,with,as,not,and,or,is,lambda,yield,True,False,None,await,async},
  sensitive=true,
  morecomment=[l]{\#},
  morestring=[b]',
  morestring=[b]"
}
\lstdefinelanguage{sql}{
  morekeywords={SELECT,FROM,WHERE,INSERT,INTO,VALUES,UPDATE,SET,DELETE,CREATE,TABLE,PRIMARY,KEY,FOREIGN,REFERENCES,UNIQUE,NOT,NULL,INDEX,ON,ALTER,DROP,CONSTRAINT},
  sensitive=false,
  morecomment=[l]{--}
}
\lstdefinelanguage{json}{
  morecomment=[l]{//},
  morestring=[b]",
  showstringspaces=false
}
\lstdefinelanguage{typescript}{
  morekeywords={import,export,from,const,let,var,function,return,if,else,for,of,in,while,switch,case,break,async,await,try,catch,finally,throw,new,class,this,extends,implements,interface,type,enum,readonly,public,private,protected,static,get,set,true,false,null,undefined,void,string,number,boolean,any,unknown,never,typeof,keyof,as},
  sensitive=true,
  morecomment=[l]{//},
  morecomment=[s]{/*}{*/},
  morestring=[b]',
  morestring=[b]",
  morestring=[b]`
}
\lstalias{javascript}{typescript}
\lstset{
  backgroundcolor=\color{codebg},
  basicstyle=\ttfamily\footnotesize,
  frame=single,
  rulecolor=\color{codeframeclr},
  frameround=tttt,
  breaklines=true,
  tabsize=2,
  showstringspaces=false,
  upquote=true,
  columns=fullflexible
}

% ---------- Helper macros ----------
\newcommand{\codefile}[1]{\texttt{#1}}
\newcommand{\kbd}[1]{\texttt{#1}}

% ---------- Code frame environment ----------
% Automatski dodaje fragile i shrink=5, koristi se umjesto \begin{frame}
% Primjer: \begin{codeframe}{Naslov} ... \end{codeframe}
% Primjer s opcijama: \begin{codeframe}[shrink=15]{Naslov} ... \end{codeframe}
\newenvironment{codeframe}[2][]
  {\begin{frame}[fragile,shrink=5,#1]{#2}}
  {\end{frame}}

% ============================================================

\documentclass[aspectratio=149,11pt]{beamer}
\usetheme{CambridgeUS}

\usepackage{fontspec}              % radi samo sa XeLaTeX/LuaLaTeX
\setsansfont{Latin Modern Sans}    % sans-serif font
\setmonofont{Latin Modern Mono}    % monospace font za kod
\usepackage[croatian]{babel}
\usepackage{csquotes}              % \enquote{} → hrvatski navodnici „..."
\usepackage{lmodern}
\usepackage{booktabs}
\usepackage{listings}
\usepackage{xcolor}
\usepackage{hyperref}

% ---------- Autor i logo ----------
\author{mag.ing. Josip Torić (SIT UNIZD)}
\titlegraphic{\includegraphics[height=2.5cm]{../common/SIT_logo_tekst1.pdf}}

% ---------- Listings ----------
\definecolor{codebg}{RGB}{248,248,248}
\definecolor{codeframeclr}{RGB}{220,220,220}
\lstdefinelanguage{bash}{
  morekeywords={sudo,apt,brew,pip,pipx,uv,winget,python,git,flask,source,deactivate,code,gunicorn,docker,compose,heroku,render,railway,node,npm,pnpm,yarn,uvicorn,alembic,pytest,ruff,black,precommit},
  sensitive=false,
  morecomment=[l]{\#}
}
\lstdefinelanguage{python}{
  morekeywords={from,import,def,class,return,if,elif,else,for,in,while,try,except,raise,with,as,not,and,or,is,lambda,yield,True,False,None,await,async},
  sensitive=true,
  morecomment=[l]{\#},
  morestring=[b]',
  morestring=[b]"
}
\lstdefinelanguage{sql}{
  morekeywords={SELECT,FROM,WHERE,INSERT,INTO,VALUES,UPDATE,SET,DELETE,CREATE,TABLE,PRIMARY,KEY,FOREIGN,REFERENCES,UNIQUE,NOT,NULL,INDEX,ON,ALTER,DROP,CONSTRAINT},
  sensitive=false,
  morecomment=[l]{--}
}
\lstdefinelanguage{json}{
  morecomment=[l]{//},
  morestring=[b]",
  showstringspaces=false
}
\lstdefinelanguage{typescript}{
  morekeywords={import,export,from,const,let,var,function,return,if,else,for,of,in,while,switch,case,break,async,await,try,catch,finally,throw,new,class,this,extends,implements,interface,type,enum,readonly,public,private,protected,static,get,set,true,false,null,undefined,void,string,number,boolean,any,unknown,never,typeof,keyof,as},
  sensitive=true,
  morecomment=[l]{//},
  morecomment=[s]{/*}{*/},
  morestring=[b]',
  morestring=[b]",
  morestring=[b]`
}
\lstalias{javascript}{typescript}
\lstset{
  backgroundcolor=\color{codebg},
  basicstyle=\ttfamily\footnotesize,
  frame=single,
  rulecolor=\color{codeframeclr},
  frameround=tttt,
  breaklines=true,
  tabsize=2,
  showstringspaces=false,
  upquote=true,
  columns=fullflexible
}

% ---------- Helper macros ----------
\newcommand{\codefile}[1]{\texttt{#1}}
\newcommand{\kbd}[1]{\texttt{#1}}

% ---------- Code frame environment ----------
% Automatski dodaje fragile i shrink=5, koristi se umjesto \begin{frame}
% Primjer: \begin{codeframe}{Naslov} ... \end{codeframe}
% Primjer s opcijama: \begin{codeframe}[shrink=15]{Naslov} ... \end{codeframe}
\newenvironment{codeframe}[2][]
  {\begin{frame}[fragile,shrink=5,#1]{#2}}
  {\end{frame}}

% ============================================================

\documentclass[aspectratio=149,11pt]{beamer}
\usetheme{CambridgeUS}

\usepackage{fontspec}              % radi samo sa XeLaTeX/LuaLaTeX
\setsansfont{Latin Modern Sans}    % sans-serif font
\setmonofont{Latin Modern Mono}    % monospace font za kod
\usepackage[croatian]{babel}
\usepackage{csquotes}              % \enquote{} → hrvatski navodnici „..."
\usepackage{lmodern}
\usepackage{booktabs}
\usepackage{listings}
\usepackage{xcolor}
\usepackage{hyperref}

% ---------- Autor i logo ----------
\author{mag.ing. Josip Torić (SIT UNIZD)}
\titlegraphic{\includegraphics[height=2.5cm]{../common/SIT_logo_tekst1.pdf}}

% ---------- Listings ----------
\definecolor{codebg}{RGB}{248,248,248}
\definecolor{codeframeclr}{RGB}{220,220,220}
\lstdefinelanguage{bash}{
  morekeywords={sudo,apt,brew,pip,pipx,uv,winget,python,git,flask,source,deactivate,code,gunicorn,docker,compose,heroku,render,railway,node,npm,pnpm,yarn,uvicorn,alembic,pytest,ruff,black,precommit},
  sensitive=false,
  morecomment=[l]{\#}
}
\lstdefinelanguage{python}{
  morekeywords={from,import,def,class,return,if,elif,else,for,in,while,try,except,raise,with,as,not,and,or,is,lambda,yield,True,False,None,await,async},
  sensitive=true,
  morecomment=[l]{\#},
  morestring=[b]',
  morestring=[b]"
}
\lstdefinelanguage{sql}{
  morekeywords={SELECT,FROM,WHERE,INSERT,INTO,VALUES,UPDATE,SET,DELETE,CREATE,TABLE,PRIMARY,KEY,FOREIGN,REFERENCES,UNIQUE,NOT,NULL,INDEX,ON,ALTER,DROP,CONSTRAINT},
  sensitive=false,
  morecomment=[l]{--}
}
\lstdefinelanguage{json}{
  morecomment=[l]{//},
  morestring=[b]",
  showstringspaces=false
}
\lstdefinelanguage{typescript}{
  morekeywords={import,export,from,const,let,var,function,return,if,else,for,of,in,while,switch,case,break,async,await,try,catch,finally,throw,new,class,this,extends,implements,interface,type,enum,readonly,public,private,protected,static,get,set,true,false,null,undefined,void,string,number,boolean,any,unknown,never,typeof,keyof,as},
  sensitive=true,
  morecomment=[l]{//},
  morecomment=[s]{/*}{*/},
  morestring=[b]',
  morestring=[b]",
  morestring=[b]`
}
\lstalias{javascript}{typescript}
\lstset{
  backgroundcolor=\color{codebg},
  basicstyle=\ttfamily\footnotesize,
  frame=single,
  rulecolor=\color{codeframeclr},
  frameround=tttt,
  breaklines=true,
  tabsize=2,
  showstringspaces=false,
  upquote=true,
  columns=fullflexible
}

% ---------- Helper macros ----------
\newcommand{\codefile}[1]{\texttt{#1}}
\newcommand{\kbd}[1]{\texttt{#1}}

% ---------- Code frame environment ----------
% Automatski dodaje fragile i shrink=5, koristi se umjesto \begin{frame}
% Primjer: \begin{codeframe}{Naslov} ... \end{codeframe}
% Primjer s opcijama: \begin{codeframe}[shrink=15]{Naslov} ... \end{codeframe}
\newenvironment{codeframe}[2][]
  {\begin{frame}[fragile,shrink=5,#1]{#2}}
  {\end{frame}}


\title[Predavanje 2: Modeliranje podataka]{Programiranje za web \\ Predavanje 2: Modeliranje podataka, ORM i migracije}
\date{15.10.2026.}

\begin{document}

\frame{\titlepage}

% ------------------------------------------------------------
\begin{frame}{Plan predavanja}
\begin{enumerate}
  \item Baze podataka: vrste, struktura, povezivanje zapisa
  \item Ograničenja i transakcije
  \item Objektno-relacijsko preslikavanje
  \item Migracije i početni podatci
\end{enumerate}

\vspace{0.6em}
Na kraju predavanja baza ima tablice \texttt{clubs} i \texttt{users}, migracija je zapisana u repozitoriju, a početni podatci se stvaraju jednom naredbom.
\end{frame}

% ------------------------------------------------------------
\begin{frame}{Ishodi učenja}
Nakon ovog predavanja student može:

\begin{itemize}
  \item razdvojiti pojmove domene u tablice i povezati ih stranim ključem
  \item obrazložiti izbor relacijske baze i izbor primarnog ključa
  \item objasniti zašto ograničenje stoji u shemi baze, a ne samo u kodu
  \item napisati model u SQLAlchemyju i razlikovati ga od sheme API-ja
  \item promijeniti shemu migracijom i objasniti doseg naredbe \texttt{downgrade}
\end{itemize}
\end{frame}

% ------------------------------------------------------------
\begin{frame}{Domena primjera}
Sustav prijava klubova na powerlifting natjecanje. Danas su potrebna samo dva pojma:

\begin{itemize}
  \item \textbf{Klub} ima naziv i grad.
  \item \textbf{Korisnik} prijavljuje se u sustav i ima ulogu. Organizator ne pripada nijednom klubu, a svaki klub ima jedan račun koji pripada samo njemu.
\end{itemize}

\vspace{0.4em}
Natjecatelji, natjecanja i prijave dolaze poslije. Prvo se postavlja temelj na kojem počivaju prijava u sustav i pravo pristupa.
\end{frame}

% ============================================================
\sekcija{Baze podataka}{Pohrana izvan aplikacije, vrste baza, struktura relacijske baze, povezivanje zapisa i ključevi.}

% ------------------------------------------------------------
\begin{frame}[fragile]{Pohrana izvan aplikacije}
Ruta iz prethodnog predavanja vraćala je popis zapisan u kodu:

\begin{lstlisting}[language=python]
@router.get("/competitions")
def list_competitions():
    return [{"id": 7, "naziv": "Zimski kup"}]
\end{lstlisting}

\vspace{0.4em}
Takav popis postoji samo u memoriji procesa, iz čega slijedi troje:

\begin{itemize}
  \item novi zapis traži izmjenu koda i ponovno pokretanje
  \item ponovno pokretanje uklanja sve što je nastalo u radu
  \item dva istodobno pokrenuta primjerka aplikacije imaju različit sadržaj
\end{itemize}

\vspace{0.3em}
Podatci zato stoje u zasebnom sustavu, koji ih pohranjuje na disk i poslužuje svim primjercima aplikacije. Takav se sustav naziva sustav za upravljanje bazom podataka.
\end{frame}

% ------------------------------------------------------------
\begin{frame}{Relacijske i nerelacijske baze}
Sustavi za upravljanje bazama dijele se prema tome kako organiziraju zapise.

\vspace{0.7em}
\begin{center}
{\small
\setlength{\tabcolsep}{8pt}
\renewcommand{\arraystretch}{1.4}
\begin{tabular}{@{}lll@{}}
  \toprule
  \textbf{Vrsta} & \textbf{Organizacija zapisa} & \textbf{Primjer sustava} \\
  \midrule
  relacijska         & tablice sa zadanom shemom  & PostgreSQL \\
  dokumentna         & dokumenti bez zadane sheme & MongoDB \\
  ključ i vrijednost & vrijednost po ključu       & Redis \\
  \bottomrule
\end{tabular}
}
\end{center}

\vspace{0.6em}
U ovom kolegiju koristi se relacijska baza jer domena ima nepromjenjiv oblik zapisa, veze među njima i pravila koja se moraju provoditi nad svakim upisom. Dokumentne baze bolje podnose zapise različita oblika, a pravila o vezama ostavljaju aplikaciji.
\end{frame}

% ------------------------------------------------------------
\begin{frame}{Struktura relacijske baze}
Relacijska baza podatke drži u tablicama. Preslikavanje pojmova domene je izravno:

\vspace{0.7em}
\begin{center}
{\small
\setlength{\tabcolsep}{8pt}
\renewcommand{\arraystretch}{1.4}
\begin{tabular}{@{}ll@{}}
  \toprule
  \textbf{Pojam domene} & \textbf{U bazi} \\
  \midrule
  vrsta zapisa      & tablica \texttt{clubs} \\
  jedan zapis       & redak \\
  svojstvo zapisa   & stupac \texttt{city} \\
  \bottomrule
\end{tabular}
}
\end{center}

\vspace{0.6em}
Tablica se izvodi iz pojma koji u domeni postoji i traje, a ne iz obrasca na zaslonu. Obrazac se mijenja češće nego što se mijenja to da klub ima naziv i grad.
\end{frame}

% ------------------------------------------------------------
\begin{frame}{Povezivanje zapisa}
Korisnik pripada klubu. Tu vezu treba zapisati tako da je baza može provjeriti.

\vspace{0.5em}
\begin{center}
\begin{tikzpicture}[
  font=\small,
  tbl/.style={draw=ink, thick, rounded corners, fill=white, inner sep=2.5mm, text=textclr},
  lbl/.style={font=\scriptsize, text=mutedclr}
]
  \node[tbl] (clubs) {%
    \begin{tabular}{@{}ll@{}}
      \multicolumn{2}{@{}l@{}}{\textbf{clubs}} \\[2pt]
      \texttt{id} & \texttt{name} \\
      \midrule
      \texttt{1} & Behemot \\
      \texttt{2} & Sokol \\
    \end{tabular}};

  \node[tbl, right=22mm of clubs] (users) {%
    \begin{tabular}{@{}lll@{}}
      \multicolumn{3}{@{}l@{}}{\textbf{users}} \\[2pt]
      \texttt{id} & \texttt{email} & \texttt{club\_id} \\
      \midrule
      \texttt{7} & ana@\dots  & \texttt{1} \\
      \texttt{8} & ivan@\dots & \texttt{1} \\
    \end{tabular}};

  \draw[-{Latex}, ink2, thick] (users.west) -- node[lbl, above] {pokazuje na} (clubs.east);
\end{tikzpicture}
\end{center}

\vspace{0.3em}
Stupac \texttt{club\_id} sadrži vrijednost stupca \texttt{id} iz tablice \texttt{clubs} i naziva se strani ključ. Baza odbija upis retka čiji \texttt{club\_id} ne postoji, kao i brisanje kluba na koji se neki redak poziva.
\end{frame}

% ------------------------------------------------------------
\begin{frame}{Primarni ključ}
Da bi jedan redak mogao pokazivati na drugi, redak treba oznaku koja ga jednoznačno određuje. Ta se oznaka naziva primarni ključ.

\vspace{0.4em}
\begin{center}
\begin{tikzpicture}[
  font=\small,
  tbl/.style={draw=ink, thick, rounded corners, fill=white, inner sep=2.5mm, text=textclr},
  lbl/.style={font=\scriptsize, text=ink, align=center}
]
  \node[tbl] (clubs) {%
    \begin{tabular}{@{}ll@{}}
      \multicolumn{2}{@{}l@{}}{\textbf{clubs}} \\[2pt]
      \texttt{id} & \texttt{name} \\
      \midrule
      \texttt{1} & Behemot \\
      \texttt{2} & Sokol \\
    \end{tabular}};

  \node[tbl, right=24mm of clubs] (users) {%
    \begin{tabular}{@{}lll@{}}
      \multicolumn{3}{@{}l@{}}{\textbf{users}} \\[2pt]
      \texttt{id} & \texttt{email} & \texttt{club\_id} \\
      \midrule
      \texttt{7} & ana@\dots  & \texttt{1} \\
      \texttt{8} & ivan@\dots & \texttt{1} \\
    \end{tabular}};

  \draw[-{Latex}, ink2, thick] (users.west) -- (clubs.east);

  \node[lbl, below=2mm of clubs] {primarni ključ \\ \texttt{clubs.id}};
  \node[lbl, below=2mm of users] {strani ključ \\ \texttt{users.club\_id}};
\end{tikzpicture}
\end{center}

\vspace{0.2em}
Veza je time opisana dvama ključevima: primarnim u tablici na koju se pokazuje i stranim u tablici koja pokazuje.
\end{frame}

% ------------------------------------------------------------
\begin{frame}{Izbor primarnog ključa}
Za primarni ključ postoje dva izbora:

\begin{itemize}
  \item \textbf{podatak iz domene}, primjerice adresa e-pošte korisnika
  \item \textbf{vrijednost bez značenja}, stupac \texttt{id} koji baza sama dodjeljuje
\end{itemize}

\vspace{0.5em}
Podatak iz domene podložan je promjeni. Promijeni li korisnik adresu e-pošte, a na njegov se redak poziva pet drugih tablica, sve te vrijednosti treba promijeniti u istoj transakciji; propust ostavlja zapise koji pokazuju na redak kojega nema.

\vspace{0.4em}
Vrijednost bez značenja nema razloga za promjenu, pa taj slučaj ne nastaje. Jedinstvenost podatka iz domene i dalje se čuva, ali ograničenjem \texttt{UNIQUE}, a ne time što je taj podatak ključ.
\end{frame}

% ------------------------------------------------------------
\begin{frame}{Operacije nad bazom}
Naredbe koje baza prima dijele se u tri skupine prema tome na što se odnose:

\vspace{0.6em}
\begin{center}
{\small
\setlength{\tabcolsep}{8pt}
\renewcommand{\arraystretch}{1.4}
\begin{tabular}{@{}lll@{}}
  \toprule
  \textbf{Skupina} & \textbf{Naredbe} & \textbf{Predmet} \\
  \midrule
  definicija sheme     & \texttt{CREATE}, \texttt{ALTER}, \texttt{DROP}   & oblik tablica \\
  rad s podatcima      & \texttt{SELECT}, \texttt{INSERT},                & sadržaj tablica \\
                       & \texttt{UPDATE}, \texttt{DELETE}                 & \\
  upravljanje transakcijom & \texttt{BEGIN}, \texttt{COMMIT}, \texttt{ROLLBACK} & početak i kraj niza izmjena \\
  \bottomrule
\end{tabular}
}
\end{center}

\vspace{0.5em}
Aplikacija u radu koristi drugu i treću skupinu. Prva se izvodi pri postavljanju baze i pri svakoj kasnijoj promjeni oblika podataka, a kako se te promjene vode, tema je četvrte cjeline.
\end{frame}

% ------------------------------------------------------------
\begin{frame}{Vrijeme kao podatak}
Rokovi, trenutci prijave i vremena izmjene zapisuju se kao datum i vrijeme. Za to postoje dva tipa stupca, i razlika među njima nije formalna:

\vspace{0.6em}
\begin{center}
{\small
\setlength{\tabcolsep}{8pt}
\renewcommand{\arraystretch}{1.4}
\begin{tabular}{@{}lll@{}}
  \toprule
  \textbf{Tip stupca} & \textbf{Što je zapisano} & \textbf{Primjer vrijednosti} \\
  \midrule
  \texttt{timestamp}   & datum i vrijeme, bez zone & \texttt{2027-02-14 23:59:00} \\
  \texttt{timestamptz} & trenutak na vremenskoj osi & \texttt{2027-02-14 23:59:00+01} \\
  \bottomrule
\end{tabular}
}
\end{center}

\vspace{0.6em}
Prva vrijednost ne određuje trenutak. Ponoć u Zadru i ponoć u Londonu zapisuju se jednako, a razlikuju se za sat. Druga vrijednost nosi odmak od koordiniranoga svjetskog vremena (UTC), pa je trenutak jednoznačan bez obzira na to tko zapis čita.
\end{frame}

% ------------------------------------------------------------
\begin{frame}{Tko provodi pravilo o zoni}
PostgreSQL vrijednost u stupcu \texttt{timestamptz} pohranjuje kao trenutak u UTC-u i pri čitanju je prikazuje u zoni koju traži onaj koji čita. Pretvorba u lokalno vrijeme time je posao prikaza, a ne pohrane.

\vspace{0.5em}
Jedno se pritom ne događa samo od sebe: \textbf{baza ne odbija vrijednost bez oznake zone}. Primi li je, tumači je u zoni koju ima poslužitelj, i upis prolazi. Pogreška se ne vidi odmah nego tek u vrijednosti koja je za sat ili dva pomaknuta.

\vspace{0.5em}
Zahtjev da vrijednost nosi oznaku zone provodi se u aplikaciji, shemom na ulazu, jednako kao i svako drugo pravilo o obliku podataka.

\vspace{0.4em}
Posljedica je konkretna: rok koji istječe u ponoć po lokalnom vremenu i rok koji istječe u ponoć po UTC-u razlikuju se za sat ili dva, a razlika je vidljiva samo onima koji zahtjev pošalju u tom razmaku.
\end{frame}

% ============================================================
\sekcija{Ograničenja i transakcije}{Uvjeti nad podatcima, istodobni zahtjevi, atomarnost, indeksi.}

% ------------------------------------------------------------
\begin{frame}{Ograničenja}
Ograničenje je uvjet zapisan u shemi tablice. Sustav ga provjerava pri svakom upisu i svakoj izmjeni.

\vspace{0.7em}
\begin{center}
{\small
\setlength{\tabcolsep}{8pt}
\renewcommand{\arraystretch}{1.4}
\begin{tabular}{@{}lll@{}}
  \toprule
  \textbf{Ograničenje} & \textbf{Uvjet} & \textbf{Primjer} \\
  \midrule
  \texttt{NOT NULL}    & vrijednost postoji              & naziv kluba \\
  \texttt{UNIQUE}      & nema dvaju jednakih             & adresa e-pošte \\
  \texttt{FOREIGN KEY} & vrijednost postoji u drugoj tablici & \texttt{club\_id} \\
  \texttt{CHECK}       & vrijednost zadovoljava izraz    & masa veća od nule \\
  \bottomrule
\end{tabular}
}
\end{center}

\vspace{0.6em}
Uvjet se provjerava neovisno o izvoru izmjene: pri upisu iz aplikacije, pri ručnom upisu iz konzole i pri uvozu iz datoteke.
\end{frame}

% ------------------------------------------------------------
\begin{frame}{Istodobni zahtjevi i jedinstvenost}
Uobičajen redoslijed u kodu je: provjeri postoji li već korisnik s tom adresom, pa upiši ako ne postoji. Dva istodobna zahtjeva tada mogu proći oba:

\vspace{0.5em}
{\small
\setlength{\tabcolsep}{7pt}
\renewcommand{\arraystretch}{1.3}
\begin{tabular}{@{}lL{4.7cm}L{4.7cm}@{}}
  \toprule
  \textbf{Trenutak} & \textbf{Zahtjev A} & \textbf{Zahtjev B} \\
  \midrule
  1 & provjera: adresa je slobodna & \\
  2 &                              & provjera: adresa je slobodna \\
  3 & upis prolazi                 & \\
  4 &                              & upis prolazi \\
  \bottomrule
\end{tabular}
}

\vspace{0.5em}
Nastaju dva retka s istom adresom, a prijava u sustav više nije jednoznačna. Nijedan zahtjev nije pogrešno napisan: razmak između provjere i upisa dovoljan je da se stanje promijeni.

\vspace{0.3em}
Uz ograničenje \texttt{UNIQUE} četvrti upis ne prolazi, a aplikacija tu pogrešku pretvara u odgovor \texttt{409}.
\end{frame}

% ------------------------------------------------------------
\begin{frame}{Nepotpun niz izmjena}
Početno punjenje baze upisuje tri retka: klub, račun organizatora i račun toga kluba.

\vspace{0.5em}
Ne prođe li treći upis, a prva dva su već potvrđena, u bazi ostaju klub i račun organizatora, bez računa kluba. To stanje nije predvidio ni kod koji upisuje ni kod koji čita: postoji klub kojim se nitko ne može služiti, a ponovno pokretanje punjenja zatječe polovicu posla obavljenom.

\vspace{0.5em}
Isto vrijedi za svaku operaciju koja mijenja više redaka. Prijava natjecatelja upisuje prijavu i umanjuje broj slobodnih mjesta; prođe li samo prvo, broj mjesta više ne odgovara broju prijava.
\end{frame}

% ------------------------------------------------------------
\begin{frame}{Transakcija}
Transakcija je niz operacija koje sustav izvodi kao jednu: ili sve budu potvrđene, ili nijedna ne ostavi traga. To se svojstvo naziva \textbf{atomarnost}.

\begin{itemize}
  \item Niz počinje naredbom \texttt{BEGIN}, a završava potvrdom (\texttt{COMMIT}) ili poništavanjem (\texttt{ROLLBACK}).
  \item Do potvrde promjene su vidljive samo unutar te transakcije.
  \item Pri pogrešci sustav vraća stanje na ono prije naredbe \texttt{BEGIN}, bez obzira na to koliko je upisa dotad prošlo.
\end{itemize}

\vspace{0.5em}
Time nepotpun niz izmjena s prethodnog slajda ne može nastati: ili su sva tri retka u bazi, ili nijedan.
\end{frame}

% ------------------------------------------------------------
\begin{frame}{Indeks}
Bez indeksa sustav pri pretrazi po stupcu čita svaki redak tablice i uspoređuje vrijednost. Indeks je zasebna struktura koja traženi redak pronalazi bez tog čitanja.

\vspace{0.6em}
\begin{center}
\begin{tikzpicture}[
  font=\scriptsize,
  redak/.style={draw=lineclr, fill=white, minimum width=16mm, minimum height=4.4mm, inner sep=0pt},
  pogodak/.style={draw=ink, thick, fill=tint, minimum width=16mm, minimum height=4.4mm, inner sep=0pt},
  nas/.style={font=\small, text=textclr}
]
  \node[nas] at (0,1.0) {bez indeksa};
  \foreach \i in {1,...,5} {
    \node[redak] (a\i) at (0, 0.5-0.55*\i) {};
    \draw[-{Latex}, ink2] (-1.7, 0.5-0.55*\i) -- (a\i.west);
  }
  \node[pogodak] at (0, 0.5-0.55*3) {};

  \node[nas] at (5.2,1.0) {s indeksom};
  \foreach \i in {1,...,5} { \node[redak] (b\i) at (5.2, 0.5-0.55*\i) {}; }
  \node[pogodak] (cilj) at (5.2, 0.5-0.55*3) {};
  \node[draw=ink, thick, rounded corners, fill=white, text=textclr] (idx) at (3.1, 0.5-0.55*3) {indeks};
  \draw[-{Latex}, ink2, thick] (idx.east) -- (cilj.west);
\end{tikzpicture}
\end{center}

\vspace{0.3em}
Primarni ključ i ograničenje \texttt{UNIQUE} indeks stvaraju sami, jer bez njega provjera jedinstvenosti ne bi bila izvediva u razumnom vremenu.
\end{frame}
% ============================================================
\sekcija{Objektno-relacijsko preslikavanje}{Uloga ORM-a, modeli, odnos modela i sheme, veza s bazom i sesija.}

% ------------------------------------------------------------
\begin{frame}[fragile]{Uloga ORM-a}
Baza prima SQL. Aplikacija radi s objektima. ORM (engl. \emph{object-relational mapping}) preslikava jedno u drugo: klasu u tablicu, primjerak klase u redak, svojstvo u stupac.

\begin{lstlisting}[language=sql]
SELECT id, name, city FROM clubs WHERE city = 'Zadar';
\end{lstlisting}

\begin{lstlisting}[language=python]
result = await session.execute(
    select(Club).where(Club.city == "Zadar")
)
clubs = result.scalars().all()
\end{lstlisting}

\vspace{0.3em}
Drugi zapis proizvodi prvi. ORM ne uklanja potrebu za poznavanjem SQL-a: kad je upit spor ili vraća neočekivan rezultat, odgovor se traži u SQL-u koji je nastao.
\end{frame}

% ------------------------------------------------------------
\begin{frame}[fragile]{Model: \codefile{Club}}
\begin{lstlisting}[language=python]
# app/models/club.py

class Club(Base):
    __tablename__ = "clubs"

    id: Mapped[int] = mapped_column(primary_key=True)
    name: Mapped[str] = mapped_column(String(120), unique=True, nullable=False)
    city: Mapped[str] = mapped_column(String(80), nullable=False)

    users: Mapped[list[User]] = relationship(back_populates="club")
\end{lstlisting}

\vspace{0.3em}
Prva tri retka opisuju stupce i ograničenja koja ulaze u shemu tablice. Posljednji redak ne stvara stupac \texttt{users}: on ORM-u govori da su korisnici toga kluba oni retci tablice \texttt{users} kojima \texttt{club\_id} odgovara ovom \texttt{id}. Zahvaljujući tomu se umjesto pisanja upita čita \texttt{club.users}, a upit sastavlja ORM.
\end{frame}

% ------------------------------------------------------------
\begin{frame}[fragile,shrink=12]{Model: \codefile{User}}
\begin{lstlisting}[language=python]
# app/models/user.py

class User(Base):
    __tablename__ = "users"

    id: Mapped[int] = mapped_column(primary_key=True)
    email: Mapped[str] = mapped_column(String(255), unique=True, nullable=False)
    password_hash: Mapped[str] = mapped_column(String(255), nullable=False)
    role: Mapped[str] = mapped_column(String(20), nullable=False)

    club_id: Mapped[Optional[int]] = mapped_column(
        ForeignKey("clubs.id"), nullable=True
    )
    club: Mapped[Optional[Club]] = relationship(back_populates="users")
\end{lstlisting}

\vspace{0.2em}
\texttt{club\_id} dopušta izostanak vrijednosti jer organizator ne pripada klubu. Zapisuje se hash lozinke, nikada sama lozinka.
\end{frame}

% ------------------------------------------------------------
\begin{frame}{Model i shema}
Dva su opisa istog pojma, s različitim zadaćama:

\begin{itemize}
  \item \textbf{model} u \codefile{models}: kako zapis izgleda u bazi
  \item \textbf{shema} u \codefile{schemas}: kako zapis putuje kroz API
\end{itemize}

\vspace{0.5em}
Razdvojeni su jer se ne poklapaju:

\begin{itemize}
  \item \texttt{password\_hash} postoji u bazi, a ne izlazi ni u jednom odgovoru
  \item pri stvaranju korisnika lozinka ulazi, pri čitanju ne izlazi, pa ulazna i izlazna shema nisu iste
  \item \texttt{id} dodjeljuje baza, pa ga klijent ne šalje, ali ga dobiva natrag
\end{itemize}

\vspace{0.4em}
Da su model i shema jedno te isto, dodavanje stupca za internu uporabu, primjerice broja neuspjelih prijava, istoga bi trena promijenilo sadržaj svakog odgovora prema klijentu.
\end{frame}

% ------------------------------------------------------------
\begin{frame}{Skup dopuštenih vrijednosti}
Stupac \texttt{role} u modelu je niz znakova, a u domeni ima samo dvije dopuštene vrijednosti, \texttt{admin} i \texttt{club}. Sama shema tablice pritom ne odbija upis vrijednosti \texttt{admn}.

\vspace{0.5em}
Posljedica nije poruka o pogrešci nego tiha promjena ponašanja. Provjera koja traži ulogu \texttt{admin} za takav redak ne prolazi, a ne prolazi ni ona koja traži \texttt{club}. Korisnik postoji, prijava u sustav uspijeva, a nijedna mu ruta nije dostupna; uzrok se ne vidi ni u jednoj poruci.

\vspace{0.5em}
Skup se može ograničiti na tri mjesta, a izbor mjesta određuje i tko je obuhvaćen i koliko stoji promjena popisa.
\end{frame}

% ------------------------------------------------------------
\begin{frame}[fragile]{Prvo mjesto: tip stupca u bazi}
\begin{lstlisting}[language=python]
# app/models/user.py
role: Mapped[str] = mapped_column(
    Enum("admin", "club", name="user_role"), nullable=False
)
\end{lstlisting}

\vspace{0.4em}
\textbf{Gdje:} u shemi tablice, opisano u sloju \codefile{models} i uneseno migracijom.

\vspace{0.2em}
\textbf{Kada se provjerava:} pri svakom upisu i izmjeni, u sustavu za upravljanje bazom.

\vspace{0.4em}
Vrijednost izvan popisa odbija se bez obzira na to odakle dolazi. Dodavanje treće uloge traži novu migraciju, jer se mijenja tip u shemi.
\end{frame}

% ------------------------------------------------------------
\begin{frame}[fragile]{Drugo mjesto: uvjet nad tablicom}
\begin{lstlisting}[language=python]
# app/models/user.py
__table_args__ = (
    CheckConstraint("role IN ('admin', 'club')", name="ck_user_role"),
)
\end{lstlisting}

\vspace{0.4em}
\textbf{Gdje:} također u shemi tablice, kao ograničenje \texttt{CHECK}.

\vspace{0.2em}
\textbf{Kada se provjerava:} pri svakom upisu i izmjeni, jednako kao tip stupca.

\vspace{0.4em}
Razlika prema prvom mjestu je u obliku, ne u dosegu: \texttt{CHECK} prima proizvoljan izraz, pa istim sredstvom pokriva i uvjete koje popis vrijednosti ne može izraziti. I ovdje promjena popisa traži migraciju.
\end{frame}

% ------------------------------------------------------------
\begin{frame}[fragile]{Treće mjesto: shema na ulazu u aplikaciju}
\begin{lstlisting}[language=python]
class UserCreate(BaseModel):
    role: Literal["admin", "club"]
\end{lstlisting}

\vspace{0.4em}
\textbf{Gdje:} u sloju \codefile{schemas}, izvan baze.

\vspace{0.2em}
\textbf{Kada se provjerava:} pri razlaganju tijela zahtjeva, prije nego što funkcija rute počne raditi.

\vspace{0.4em}
Odgovor imenuje polje koje nije prošlo, a promjena popisa ne traži migraciju. Zauzvrat, upis koji ne dolazi kroz aplikaciju, primjerice iz konzole baze, ovim uvjetom nije obuhvaćen.

\vspace{0.3em}
\begin{konvencija}[Izbor, ne pravilo]
Nijedno mjesto nije propisano. U ovom projektu popis nije zapisan nigdje: stupac \texttt{role} je \texttt{String(20)}, a komentar u modelu to obrazlaže time da je bez nabrajanja jednostavnije za početak.
\end{konvencija}
\end{frame}

% ------------------------------------------------------------
\begin{frame}[fragile]{Veza s bazom i sesija}
\begin{lstlisting}[language=python]
# app/core/database.py
engine = create_async_engine(settings.DATABASE_URL, pool_pre_ping=True)
AsyncSessionLocal = async_sessionmaker(bind=engine, class_=AsyncSession)

# app/core/deps.py
async def get_db():
    async with AsyncSessionLocal() as session:
        try:
            yield session
            await session.commit()
        except Exception:
            await session.rollback()
            raise
\end{lstlisting}

\vspace{0.2em}
\texttt{engine} drži skup otvorenih veza prema bazi. \texttt{session} je jedna cjelina rada: prati promijenjene objekte i na kraju ih zajedno potvrđuje ili poništava.
\end{frame}

% ------------------------------------------------------------
\begin{frame}{Doseg sesije}
Sesija nije samo veza prema bazi. Ona prikuplja promijenjene objekte i unutar nje teče jedna transakcija, pa je smije koristiti samo jedan zahtjev.

\begin{itemize}
  \item Dva zahtjeva nad istom sesijom dijele i transakciju: potvrda jednoga potvrđuje i nedovršene promjene drugoga.
  \item Poništavanje zbog pogreške u jednom zahtjevu poništilo bi i rad drugoga.
\end{itemize}

\vspace{0.5em}
Zato sesiju otvara ovisnost pri ulasku u rutu i zatvara je pri izlasku.
\end{frame}

% ------------------------------------------------------------
\begin{frame}{Čekanje na bazu}
Upit prema bazi traje, a proces za to vrijeme nema što raditi. Dva su načina da se to čekanje podnese:

\begin{itemize}
  \item \textbf{sinkrono}: dretva koja je poslala upit čeka odgovor i u tom razdoblju ne poslužuje druge zahtjeve; istodobnost se postiže većim brojem dretvi ili procesa
  \item \textbf{asinkrono}: dretva za vrijeme čekanja preuzima drugi zahtjev i vraća se kad odgovor stigne
\end{itemize}

\vspace{0.5em}
Izbor nije stvar ukusa nego mora biti dosljedan kroz cijeli sloj pristupa bazi: asinkroni kod traži upravljački program koji podržava takav rad. Adresa za spajanje ga i imenuje, pa \texttt{postgresql+asyncpg://} označava asinkroni, a \texttt{postgresql://} obični.

\vspace{0.3em}
U ovom se projektu radi asinkrono, pa su i funkcije koje dodiruju bazu označene s \texttt{async}, a njihovi pozivi s \texttt{await}.
\end{frame}

% ============================================================
\sekcija{Migracije i početni podatci}{Razilaženje modela i baze, sadržaj migracije, naredbe, granice automatskog stvaranja, punjenje baze.}

% ------------------------------------------------------------
\begin{frame}{Razilaženje modela i baze}
Dodavanje stupca u model mijenja Python kod. Tablica u bazi ostaje kakva je bila, pa prvi upit koji taj stupac spominje pada.

\vspace{0.5em}
Brisanje i ponovno stvaranje tablica rješava to samo dok je baza prazna. Čim u njoj postoje podatci, potreban je zapis koji staru shemu prevodi u novu bez gubitka sadržaja.

\vspace{0.5em}
Migracija je takav zapis: datoteka u repozitoriju koja opisuje jednu promjenu sheme. Iz toga slijedi troje:

\begin{itemize}
  \item svaki član tima izvodi iste promjene istim redoslijedom
  \item produkcijska baza prolazi kroz iste korake kao razvojna
  \item promjena sheme vidljiva je u povijesti, uz kod koji ju je tražio
\end{itemize}
\end{frame}

% ------------------------------------------------------------
\begin{frame}[fragile,shrink=15]{Sadržaj migracije}
\begin{lstlisting}[language=python]
# alembic/versions/001_init_clubs_and_users.py
revision = "001"
down_revision = None

def upgrade() -> None:
    op.create_table(
        "clubs",
        sa.Column("id", sa.Integer(), primary_key=True),
        sa.Column("name", sa.String(length=120), nullable=False),
        sa.Column("city", sa.String(length=80), nullable=False),
        sa.UniqueConstraint("name"),
    )

def downgrade() -> None:
    op.drop_table("clubs")
\end{lstlisting}

\vspace{0.2em}
\texttt{revision} imenuje ovu promjenu, \texttt{down\_revision} onu prije nje; iz toga slijedi redoslijed. Baza pamti do koje je migracije došla, pa svaka prolazi točno jednom.
\end{frame}

% ------------------------------------------------------------
\begin{frame}[fragile]{Naredbe}
\begin{lstlisting}[language=bash]
# create a migration from the difference between models and database
alembic revision --autogenerate -m "init clubs and users"

# apply everything not yet applied
alembic upgrade head

# undo the last migration
alembic downgrade -1

# which migration is the database on
alembic current
\end{lstlisting}

\vspace{0.4em}
\texttt{upgrade head} primjenjuje sve migracije kroz koje baza još nije prošla, redoslijedom koji slijedi iz njihovih oznaka.
\end{frame}

% ------------------------------------------------------------
\begin{frame}{Granice automatskog stvaranja}
\texttt{autogenerate} uspoređuje modele s trenutnim stanjem baze i zapisuje razliku. Za neke promjene ta usporedba nije dostatna:

\begin{itemize}
  \item Promjena naziva stupca vidi se kao nestanak jednog i pojava drugog, pa nastaju brisanje i stvaranje.
  \item Promjena tipa stupca traži i pravilo po kojem se postojeće vrijednosti pretvaraju, a ono iz usporedbe ne slijedi.
  \item Novi stupac s oznakom \texttt{NOT NULL} nad tablicom koja već ima retke ne može se dodati bez zadane vrijednosti, jer bi postojeći retci ostali bez vrijednosti. Nad praznom tablicom isti upis prolazi.
\end{itemize}

\vspace{0.5em}
Prva dva slučaja ne javljaju pogrešku. Migracija se izvede do kraja i \texttt{upgrade} prođe uredno; nedostaju samo podatci.

\vspace{0.3em}
Doseg naredbe \texttt{downgrade} pritom je ograničen: brisanje stupca ne može se poništiti jer podatci koje je sadržavao više ne postoje.
\end{frame}

% ------------------------------------------------------------
\begin{frame}[fragile]{Čitanje migracije prije pokretanja}
Nakon preimenovanja stupca \texttt{email} u \texttt{username}, \texttt{autogenerate} zapisuje ovo:

\begin{lstlisting}[language=python]
def upgrade() -> None:
    op.drop_column("users", "email")
    op.add_column("users", sa.Column("username", sa.String(255)))
\end{lstlisting}

\vspace{0.3em}
Znak za uzbunu je par \texttt{drop\_column} i \texttt{add\_column} nad istom tablicom u istoj migraciji: to je preimenovanje koje je izgubilo sadržaj. Ispravak je jedan redak:

\begin{lstlisting}[language=python]
def upgrade() -> None:
    op.alter_column("users", "email", new_column_name="username")
\end{lstlisting}

\vspace{0.3em}
Zato se svaka nastala migracija pročita prije pokretanja. U datoteci je ispravak trenutačan; nakon naredbe \texttt{upgrade} nad produkcijskom bazom sadržaja stupca više nema.
\end{frame}

% ------------------------------------------------------------
\begin{frame}[fragile]{Dodavanje indeksa}
Indeks se traži uz stupac u modelu, a u shemu ulazi migracijom:

\begin{lstlisting}[language=python]
# in the model
club_id: Mapped[int] = mapped_column(ForeignKey("clubs.id"), index=True)

# in the migration
op.create_index("ix_users_club_id", "users", ["club_id"])
\end{lstlisting}

\vspace{0.4em}
Indeks ubrzava pretragu, spajanje tablica i sortiranje po tom stupcu. Zauzima prostor i usporava upis, jer se pri svakoj izmjeni i sam mijenja.

\vspace{0.3em}
\begin{konvencija}
Indeks se dodaje na stupce po kojima se često filtrira, tipično na strane ključeve. Indeksiranje svakog stupca unaprijed plaća se pri svakom upisu, a rijetko donosi korist.
\end{konvencija}
\end{frame}


% ------------------------------------------------------------
\begin{frame}[fragile,shrink=12]{Početni podatci}
Za isprobavanje aplikacije potrebni su organizator, jedan klub i račun toga kluba. Migracija ih ne stvara jer mijenja oblik tablica, a ne njihov sadržaj; postupak koji ih upisuje naziva se punjenje početnim podatcima (engl. \emph{seeding}).

\begin{lstlisting}[language=python]
# app/seed.py
result = await session.execute(
    select(Club).where(Club.name == DEMO_CLUB_NAME)
)
club = result.scalar_one_or_none()
if club is None:
    club = Club(name=DEMO_CLUB_NAME, city=DEMO_CLUB_CITY)
    session.add(club)
    await session.flush()
\end{lstlisting}

\vspace{0.2em}
Provjera prije upisa čini postupak ponovljivim: drugo pokretanje ne stvara drugi klub i ne pada na ograničenju \texttt{UNIQUE}. \texttt{flush} šalje upis bazi prije potvrde transakcije, da klub dobije \texttt{id} koji korisniku treba kao strani ključ.
\end{frame}

% ------------------------------------------------------------
\begin{frame}{Sažetak}
\begin{itemize}
  \item Relacijska baza odabrana je jer domena ima nepromjenjiv oblik zapisa, veze među njima i pravila koja se provode nad svakim upisom.
  \item Tablica se izvodi iz pojma koji u domeni traje. Redak se određuje vrijednošću bez značenja, jer se podatci iz domene mijenjaju, a veze se opisuju parom primarnog i stranog ključa.
  \item Ograničenje u shemi vrijedi za svaki upis. Provjera u kodu ne obuhvaća dva istodobna zahtjeva, jer se stanje može promijeniti između provjere i upisa.
  \item ORM preslikava klasu u tablicu i sastavlja SQL, ali ga ne skriva. Model baze i shema API-ja ostaju razdvojeni.
  \item Migracija je zapis jedne promjene sheme u repozitoriju. Čita se prije pokretanja, jer \texttt{downgrade} ne vraća obrisane podatke.
\end{itemize}
\end{frame}

% ------------------------------------------------------------
\begin{frame}{Što slijedi}
\begin{itemize}
  \item \textbf{22.10.2026.} Predavanje 3, autentikacija: pohrana lozinke, prijava, token i njegova provjera pri svakom zahtjevu.
  \item \textbf{29.10.2026.} Predavanje 4, autorizacija: uloge i vlasništvo nad zapisom.
  \item \textbf{19.11.2026.} Obrana prvoga dijela projekta.
\end{itemize}

\vspace{0.5em}
Do sljedećega predavanja: pokrenuti bazu, primijeniti migracije naredbom \texttt{alembic upgrade head} i potvrditi da u tablicama postoje organizator i demo klub.
\end{frame}

\end{document}
